\chapter{Architecture et implémentation du système}

\section*{Introduction}
\addcontentsline{toc}{section}{Introduction}

Le chapitre précédent a défini le modèle multi-agents, les responsabilités, les interactions et les règles d'autorité. Le présent chapitre décrit leur mise en œuvre concrète dans les deux systèmes de la Migration Factory. Il présente l'environnement technique, la structure des dépôts, l'implémentation des agents, l'orchestrateur, le backend, la persistance, le cockpit et les principales difficultés rencontrées.

Les deux solutions ne possèdent pas le même niveau de maturité. Java/Spring Boot constitue la référence la plus avancée, avec un workflow complet, une réparation gouvernée, un assistant et un reporting. Angular est en développement actif : les agents principaux sont implémentés et l'intégration du workflow de bout en bout est en cours. Les sections suivantes signalent cette différence afin de ne pas présenter une fonctionnalité en cours comme un résultat achevé.

% ============================================================
\section{Environnement technique}
% ============================================================

\subsection{Technologies communes}

\begin{table}[H]
    \centering
    \small
    \arrayrulecolor{cgigray}
    \rowcolors{2}{cgilight!55}{white}
    \begin{tabularx}{\textwidth}{|>{\bfseries\RaggedRight\arraybackslash}p{3.8cm}|>{\RaggedRight\arraybackslash}p{4.4cm}|>{\RaggedRight\arraybackslash}X|}
        \hline
        \rowcolor{cgilight!95}
        \textbf{Couche} & \textbf{Technologies} & \textbf{Rôle} \\
        \hline
        Interface & React et Next.js & Configuration, cockpit, décisions, assistant et rapports. \\
        \hline
        Backend & Python et FastAPI & API HTTP, validation, services applicatifs et projections. \\
        \hline
        Orchestration & LangGraph & Graphe d'exécution, interruptions, reprises et routage. \\
        \hline
        Persistance & SQLite & États, événements, décisions, commandes et références d'artefacts. \\
        \hline
        Intelligence & Azure OpenAI / modèles configurés par rôle & Analyse, planification, review, réparation et assistance. \\
        \hline
        Validation & Tests Python et tests frontend & Vérification des services, contrats, API et composants. \\
        \hline
        Collaboration & GitHub, Jira, Teams et Confluence & Versionnement, backlog, communication et documentation. \\
        \hline
    \end{tabularx}
    \caption{Technologies communes aux deux systèmes}
    \label{tab:technologies-communes-ch4}
    \rowcolors{2}{white}{white}
    \arrayrulecolor{black}
\end{table}

FastAPI fournit les routes et l'injection de dépendances utilisées pour assembler les services applicatifs et les adaptateurs \cite{fastapi-dependencies}. LangGraph assure la coordination des étapes avec un état et des checkpoints, ce qui facilite les interruptions humaines et la reprise \cite{langgraph-persistence}. SQLite reste adapté à un prototype local ; le mode WAL réduit certains blocages entre lecteurs et écrivain \cite{sqlite-wal}.

\subsection{Environnement Java/Spring Boot}

La chaîne Java mobilise :

\begin{itemize}
    \item JDK 11, 17 et 21 selon le stage ;
    \item Maven pour la résolution des dépendances, la compilation et les tests ;
    \item OpenRewrite pour des transformations déterministes composables ;
    \item des profils Spring Boot décrivant les trajectoires 2.1, 2.7, 3.5 et 4.0 ;
    \item PowerShell pour plusieurs scripts locaux de lancement et de diagnostic sous Windows.
\end{itemize}

OpenRewrite applique des recettes sur des arbres sémantiques préservant le format et permet de composer plusieurs opérations pour réaliser une migration plus large \cite{openrewrite-recipes}. Le système l'encapsule dans un workflow où l'exécution, le répertoire et les preuves sont contrôlés par le backend.

\subsection{Environnement Angular}

La chaîne Angular mobilise :

\begin{itemize}
    \item Node.js et npm, sélectionnés selon la compatibilité du stage ;
    \item Angular CLI et \texttt{ng update} ;
    \item TypeScript et RxJS ;
    \item les fichiers \texttt{package.json}, \texttt{angular.json}, \texttt{tsconfig.json} et le lockfile ;
    \item des workspaces nettoyés de \texttt{node\_modules}, \texttt{dist}, caches et couvertures ;
    \item des commandes d'installation, de build et de test résolues par le backend.
\end{itemize}

La matrice officielle Angular relie chaque version à des plages compatibles de Node.js, TypeScript et RxJS \cite{angular-versions}. Cette contrainte est implémentée sous la forme d'un catalogue de profils et d'un service de résolution.

% ============================================================
\section{Structure des projets}
% ============================================================

\subsection{Dépôt Java/Spring Boot}

La structure logique du dépôt Java sépare domaine, application, infrastructure, orchestration, agents, interface et tests.

\begin{verbatim}
migration-factory-java/
|-- migration_factory/
|   |-- control_tower/
|   |   |-- adapters/fastapi/
|   |   |-- application/
|   |   |-- domain/
|   |   `-- infrastructure/sqlite/
|   |-- orchestrator/
|   `-- agents/
|       |-- analysis_agent/
|       |-- planning_agent/
|       |-- transformation_agent/
|       `-- repair_agents/
|-- web/control-tower/
|-- tests/control_tower/
|-- scripts/
`-- docs/
\end{verbatim}

Le dossier \texttt{control\_tower} centralise les cas d'utilisation et les projections du cockpit. L'orchestrateur coordonne les nœuds. Les agents restent séparés des adaptateurs HTTP et SQLite. Le frontend communique uniquement avec les routes exposées par le backend.

\subsection{Dépôt Angular}

La structure Angular suit les mêmes principes avec des services propres aux profils Node/npm et aux stages Angular.

\begin{verbatim}
angular-migration/
|-- backend/
|   |-- app/
|   |   |-- domain/
|   |   |-- services/
|   |   |-- api/
|   |   |-- infrastructure/
|   |   `-- orchestration/
|   `-- tests/
|-- frontend/
|-- goals/
|-- artifacts/
|-- scripts/
`-- tests/
\end{verbatim}

Les deux dépôts ne partagent pas leurs services d'exécution. Le partage concerne les principes, la terminologie et les règles de gouvernance. Cette indépendance permet aux catalogues et agents Angular d'évoluer sans perturber le workflow Java.

\subsection{Vue de déploiement local}

\begin{figure}[H]
    \centering
    \resizebox{0.94\textwidth}{!}{%
    \begin{tikzpicture}[
        node/.style={draw=cgiblue,rounded corners=3pt,fill=cgilight!55,minimum width=4.0cm,minimum height=1.15cm,align=center,font=\small\bfseries},
        store/.style={draw=cgigray,rounded corners=3pt,fill=white,minimum width=4.0cm,minimum height=1.15cm,align=center,font=\small\bfseries},
        arr/.style={-{Latex[length=2mm]},thick,draw=cgigray}
    ]
        \node[node] (browser) {Navigateur\\cockpit local};
        \node[node,right=1.0cm of browser] (frontend) {Serveur frontend\\Next.js};
        \node[node,right=1.0cm of frontend] (backend) {Backend local\\FastAPI};
        \node[store,below=1.1cm of backend] (db) {SQLite et artefacts};
        \node[store,left=1.0cm of db] (workspace) {Workspaces\\et outils};
        \node[store,right=1.0cm of db] (llm) {Fournisseur LLM\\externe};
        \draw[arr] (browser)--(frontend);
        \draw[arr] (frontend)--(backend);
        \draw[arr] (backend)--(db);
        \draw[arr] (backend)--(workspace);
        \draw[arr] (backend)--(llm);
    \end{tikzpicture}%
    }
    \caption{Déploiement local des prototypes}
    \label{fig:deploiement-local-ch4}
\end{figure}

% ============================================================
\section{Développement des agents}
% ============================================================

Chaque agent est décrit à l'aide d'une fiche uniforme : objectif, entrées, traitement, sorties, scénario et limitations.

\subsection{Agent d'analyse}

\begin{table}[H]
    \centering
    \small
    \begin{tabularx}{\textwidth}{|>{\bfseries\RaggedRight\arraybackslash}p{3.2cm}|>{\RaggedRight\arraybackslash}X|}
        \hline
        Objectif & Construire une représentation exploitable de l'application sans modifier la source. \\
        \hline
        Entrées & Snapshot, versions, manifestes, structure, configurations et profil d'exécution. \\
        \hline
        Traitement & Inventaire déterministe, extraction de faits, construction d'un contexte borné et enrichissement LLM. \\
        \hline
        Sorties & Rapport d'analyse structuré, risques, dépendances, preuves et références de fichiers. \\
        \hline
        Scénario & Détecter une ancienne version Spring Boot ou une incompatibilité Angular--Node.js. \\
        \hline
        Limitations & Les dépendances privées ou configurations dynamiques peuvent nécessiter une information humaine. \\
        \hline
    \end{tabularx}
    \caption{Fiche de l'agent d'analyse}
    \label{tab:agent-analyse-ch4}
\end{table}

Dans Java, l'analyse examine notamment le \texttt{pom.xml}, le code, les propriétés et les tests. Dans Angular, elle analyse \texttt{package.json}, \texttt{angular.json}, TypeScript et les profils de compatibilité.

\subsection{Reviewer d'analyse}

Le reviewer reçoit le rapport et les preuves. Il vérifie la couverture, la cohérence des versions et la présence des informations nécessaires au planificateur. Sa sortie est une décision structurée : accepter, demander une révision ou refuser. Il ne lance aucune commande et ne modifie pas le projet.

\subsection{Agent de planification}

\begin{table}[H]
    \centering
    \small
    \begin{tabularx}{\textwidth}{|>{\bfseries\RaggedRight\arraybackslash}p{3.2cm}|>{\RaggedRight\arraybackslash}X|}
        \hline
        Objectif & Décomposer la cible en stages, unités de travail et validations. \\
        \hline
        Entrées & Analyse validée, cible, catalogue de compatibilité et politiques. \\
        \hline
        Traitement & Sélection d'une trajectoire, ordonnancement des transformations et définition des critères de sortie. \\
        \hline
        Sorties & Plan versionné, dépendances, risques, commandes attendues et points de décision. \\
        \hline
        Scénario & Construire 2.1 $\rightarrow$ 2.7 $\rightarrow$ 3.5 $\rightarrow$ 4.0 ou Angular 18 $\rightarrow$ 19 $\rightarrow$ 20 $\rightarrow$ 21. \\
        \hline
        Limitations & Le plan dépend de la qualité de l'analyse et des profils disponibles. \\
        \hline
    \end{tabularx}
    \caption{Fiche de l'agent de planification}
    \label{tab:agent-planification-ch4}
\end{table}

Le reviewer du plan contrôle ensuite l'ordre, les prérequis et les validations. Un plan rejeté produit une nouvelle version ; l'ancien plan n'est pas modifié silencieusement.

\subsection{Agent de transformation}

L'agent de transformation consomme uniquement un plan approuvé et un workspace autorisé. Dans Java, il combine des transformations Maven, Java et OpenRewrite. Dans Angular, il prépare les modifications des dépendances, configurations et sources TypeScript. Sa sortie est un ensemble de changements candidats accompagné de métadonnées.

L'agent ne choisit pas librement un répertoire. Le chemin est fourni par le backend et vérifié avant écriture. Les transformations déterministes peuvent être appliquées directement par un service ; les modifications proposées par un modèle sont converties en artefacts contrôlables.

\subsection{Agent de build et de tests}

La validation repose sur un manifeste de commande construit par le backend. L'agent ou service de build ne reçoit pas une chaîne shell arbitraire. Il exécute Maven pour Java ou npm/Angular CLI pour Angular, puis capture : code de sortie, durée, journaux, résultats de tests et fichiers produits.

Un succès est déterminé par les résultats d'exécution, non par une affirmation du LLM. En cas d'échec, les preuves sont classées et transmises au workflow de réparation.

\subsection{Agents de réparation}

Le \textbf{proposeur} reçoit un contexte d'échec borné : erreurs, extraits de code, commande, profil et état du workspace. Il produit un diff structuré. Le \textbf{reviewer} vérifie la cohérence, les risques et l'applicabilité ; il peut accepter, refuser ou consolider le correctif. La version finale reste soumise au développeur.

\begin{figure}[H]
    \centering
    \resizebox{0.96\textwidth}{!}{%
    \begin{tikzpicture}[
        step/.style={draw=cgiblue,rounded corners=3pt,fill=cgilight!55,minimum width=3.0cm,minimum height=1.0cm,align=center,font=\small\bfseries},
        human/.style={draw=cgired,rounded corners=3pt,fill=white,minimum width=3.0cm,minimum height=1.0cm,align=center,font=\small\bfseries},
        arr/.style={-{Latex[length=2mm]},thick,draw=cgigray}
    ]
        \node[step] (failure) {Échec vérifié};
        \node[step,right=0.45cm of failure] (context) {Contexte borné};
        \node[step,right=0.45cm of context] (proposal) {Proposition};
        \node[step,right=0.45cm of proposal] (review) {Review};
        \node[human,right=0.45cm of review] (human) {Approbation};
        \node[step,below=1.1cm of human] (apply) {Application};
        \node[step,left=0.45cm of apply] (validate) {Build et tests};
        \draw[arr] (failure)--(context);
        \draw[arr] (context)--(proposal);
        \draw[arr] (proposal)--(review);
        \draw[arr] (review)--(human);
        \draw[arr] (human)--(apply);
        \draw[arr] (apply)--(validate);
        \draw[arr,bend left=30] (validate.north) to (context.south);
    \end{tikzpicture}%
    }
    \caption{Implémentation de la chaîne de réparation}
    \label{fig:implementation-reparation-ch4}
\end{figure}

Dans l'implémentation Java actuelle, le reviewer peut produire une version consolidée du diff. Cette décision est contrôlée par des schémas, une lignée, des empreintes et une approbation humaine. La réparation Angular complète reste à intégrer.

\subsection{Assistant de migration}

L'assistant utilise une projection de l'état courant, les derniers événements, les preuves autorisées et un historique borné de conversation. Il peut expliquer un blocage, résumer les étapes et indiquer la prochaine action. Les prompts imposent un comportement en lecture seule : l'assistant ne peut pas exécuter, approuver ou modifier des fichiers.

Sa principale limitation est la qualité des données projetées. Une information absente de l'état ou des artefacts ne doit pas être inventée. Les réponses sont donc présentées comme une assistance, non comme une preuve d'exécution.

% ============================================================
\section{Développement de l'orchestrateur}
% ============================================================

\subsection{Réception de la requête}

Le frontend transmet une intention à l'API : créer une migration, démarrer un stage, approuver, rejeter, réviser, annuler ou reprendre. Le backend valide le schéma, charge la migration et vérifie que l'action est autorisée dans l'état courant.

\subsection{Choix des agents}

Le choix n'est pas effectué par une discussion libre. Le nœud suivant dépend du stage et de l'état persistant. L'orchestrateur appelle l'analyseur après la qualification, le reviewer après l'analyse, le planificateur après validation, puis le transformateur après approbation.

\subsection{Séquencement des tâches}

\begin{figure}[H]
    \centering
    \resizebox{\textwidth}{!}{%
    \begin{tikzpicture}[
        step/.style={draw=cgiblue,rounded corners=3pt,fill=cgilight!55,minimum width=2.75cm,minimum height=1.0cm,align=center,font=\small\bfseries},
        gate/.style={draw=cgired,rounded corners=3pt,fill=white,minimum width=2.75cm,minimum height=1.0cm,align=center,font=\small\bfseries},
        arr/.style={-{Latex[length=2mm]},thick,draw=cgigray}
    ]
        \node[step] (qualify) {Qualification};
        \node[step,right=0.4cm of qualify] (analyse) {Analyse};
        \node[step,right=0.4cm of analyse] (review) {Review};
        \node[step,right=0.4cm of review] (plan) {Planification};
        \node[gate,right=0.4cm of plan] (approval) {Décision};
        \node[step,right=0.4cm of approval] (transform) {Transformation};
        \node[step,right=0.4cm of transform] (validate) {Validation};
        \draw[arr] (qualify)--(analyse);
        \draw[arr] (analyse)--(review);
        \draw[arr] (review)--(plan);
        \draw[arr] (plan)--(approval);
        \draw[arr] (approval)--(transform);
        \draw[arr] (transform)--(validate);
    \end{tikzpicture}%
    }
    \caption{Séquencement principal implémenté par l'orchestrateur}
    \label{fig:sequencement-orchestrateur-ch4}
\end{figure}

\subsection{Récupération des résultats}

Chaque nœud publie un résultat structuré et des artefacts. L'orchestrateur enregistre les références, met à jour le stage et produit un événement. Le frontend obtient ensuite une projection cohérente à partir de la base et non directement à partir de la réponse du modèle.

\subsection{Gestion des erreurs}

Les erreurs sont classées selon leur origine : configuration, environnement, agent, commande, transformation, build ou test. Une erreur de précondition bloque avant l'exécution. Une erreur de modèle produit un échec fermé ou un fallback. Une erreur de build ouvre la réparation si les règles et le nombre de tentatives le permettent.

\subsection{Production du résultat final}

Lorsque tous les stages sont validés, un service de rapport agrège versions, étapes, décisions, artefacts, résultats de build et tests, réparations et références de preuves. Le résultat reste un candidat de migration ; le système ne réalise pas de déploiement en production.

% ============================================================
\section{Backend, persistance et contrats}
% ============================================================

\subsection{API d'intégration}

Les API sont organisées autour des ressources suivantes :

\begin{itemize}
    \item configurations et prérequis ;
    \item migrations, stages et progression ;
    \item points de décision et approbations ;
    \item propositions de réparation ;
    \item artefacts, événements et rapports ;
    \item assistant et consommation des modèles.
\end{itemize}

Les requêtes utilisent des schémas typés. Les décisions sensibles incluent une empreinte attendue afin de détecter qu'un artefact a changé depuis son affichage.

\subsection{Modèle de données}

\begin{figure}[H]
    \centering
    \resizebox{0.95\textwidth}{!}{%
    \begin{tikzpicture}[
        entity/.style={draw=cgiblue,rounded corners=3pt,fill=cgilight!55,minimum width=2.8cm,minimum height=0.95cm,align=center,font=\small\bfseries},
        arr/.style={-{Latex[length=2mm]},thick,draw=cgigray}
    ]
        \node[entity] (run) {Migration};
        \node[entity,right=0.7cm of run] (stage) {Stage};
        \node[entity,right=0.7cm of stage] (command) {Commande};
        \node[entity,right=0.7cm of command] (artifact) {Artefact};
        \node[entity,below=1.0cm of run] (event) {Événement};
        \node[entity,below=1.0cm of stage] (gate) {Gate};
        \node[entity,below=1.0cm of command] (proposal) {Proposition};
        \node[entity,below=1.0cm of artifact] (report) {Rapport};
        \draw[arr] (run)--(stage);
        \draw[arr] (stage)--(command);
        \draw[arr] (command)--(artifact);
        \draw[arr] (run)--(event);
        \draw[arr] (stage)--(gate);
        \draw[arr] (gate)--(proposal);
        \draw[arr] (artifact)--(report);
    \end{tikzpicture}%
    }
    \caption{Principales entités persistées}
    \label{fig:entites-persistantes-ch4}
\end{figure}

Les événements sont ajoutés de manière séquentielle. Les artefacts stockent leur type, producteur, chemin interne et empreinte. Les commandes conservent leur état, durée, code de sortie et références de journaux.

\subsection{Workspaces et artefacts}

Chaque migration possède un répertoire dédié, subdivisé par stage et type d'artefact.

\begin{verbatim}
run_<id>/
|-- source_snapshot/
|-- stages/
|   |-- stage_01/workspace/
|   |-- stage_01/analysis/
|   |-- stage_01/planning/
|   |-- stage_01/validation/
|   `-- stage_02/...
|-- repairs/
|-- logs/
|-- evidence/
`-- reports/
\end{verbatim}

La source externe n'est pas utilisée comme répertoire d'écriture. Le candidat validé d'un stage sert de base au suivant. Les artefacts importants sont liés par SHA-256.

% ============================================================
\section{Interface utilisateur et cockpit}
% ============================================================

\subsection{Configuration}

L'écran de création collecte les chemins, versions source et cible, profils et options. Les valeurs sont validées par le backend. Dans Java, la source est fournie par chemin local. Angular utilise également un projet local ou un package préparé selon le workflow en cours d'intégration.

\subsection{Tableau de bord}

Le cockpit présente :

\begin{itemize}
    \item résumé de la migration ;
    \item pipeline et stage sélectionné ;
    \item analyse, plan et points de décision ;
    \item commandes, journaux et résultats ;
    \item événements du workflow ;
    \item propositions de réparation ;
    \item assistant et rapports.
\end{itemize}

\begin{figure}[H]
    \centering
    \resizebox{0.94\textwidth}{!}{%
    \begin{tikzpicture}[
        screen/.style={draw=cgigray,rounded corners=4pt,fill=white,minimum width=14cm,minimum height=7cm},
        nav/.style={draw=cgiblue,fill=cgiblue!10,minimum width=3.0cm,minimum height=6.4cm,align=center,font=\small},
        panel/.style={draw=cgilight,fill=cgilight!35,rounded corners=2pt,minimum width=4.6cm,minimum height=2.0cm,align=center,font=\small\bfseries}
    ]
        \node[screen] (screen) {};
        \node[nav] at ($(screen.west)+(1.7,0)$) {Navigation\\[2mm]Vue d'ensemble\\Pipeline\\Analyse\\Commandes\\Événements\\Assistant\\Rapports};
        \node[panel] at ($(screen.center)+(1.0,1.8)$) {Résumé et état courant};
        \node[panel] at ($(screen.center)+(6.0,1.8)$) {Décision ou blocage};
        \node[panel,minimum width=9.8cm] at ($(screen.center)+(3.5,-1.0)$) {Stage sélectionné, preuves, logs et prochaines actions};
        \node[font=\small\bfseries] at ($(screen.north)+(0,-0.35)$) {Cockpit de migration};
    \end{tikzpicture}%
    }
    \caption{Organisation fonctionnelle du cockpit}
    \label{fig:cockpit-ch4}
\end{figure}

Le frontend ne conserve pas un état métier parallèle. Il projette les ressources fournies par le backend et envoie des intentions explicites.

\subsection{Logs et preuves}

Les journaux de commande sont stockés comme artefacts et affichés sous forme de vues bornées. Les événements permettent de suivre l'ordre des transitions. Les preuves sensibles sont accompagnées d'une empreinte et d'une provenance.

% ============================================================
\section{Authentification et sécurité}
% ============================================================

Les prototypes ne comportent pas encore d'authentification, de comptes ni de contrôle d'accès par rôle. Le système Java limite l'API au poste local, vérifie l'origine des requêtes mutatives et exige un identifiant client statique. Ces mécanismes protègent l'usage local, mais ne constituent pas une authentification.

L'identité utilisée dans certaines traces provient du compte du système d'exploitation ou de valeurs génériques. Elle ne permet pas une attribution fiable à plusieurs utilisateurs. Une industrialisation devra ajouter :

\begin{itemize}
    \item authentification forte ;
    \item rôles développeur, reviewer et administrateur ;
    \item autorisation par action ;
    \item journal d'audit attribuable ;
    \item gestion des sessions et révocation.
\end{itemize}

Les contextes LLM sont limités et nettoyés. Les fichiers de secrets, clés et chemins interdits sont exclus. Les redactions restent fondées sur des règles et ne garantissent pas la détection de toute valeur sensible intégrée arbitrairement dans le code.

% ============================================================
\section{Principales difficultés techniques et corrections}
% ============================================================

\subsection{Application fiable des correctifs}

Les premières versions de la réparation Java ont rencontré des différences de fins de ligne, des hunks non applicables et des propositions liées à un état obsolète. Les corrections ont introduit des diffs structurés, des empreintes, une lignée de propositions, une validation d'applicabilité et une application uniquement dans le sandbox.

\subsection{Autorité des commandes}

Plusieurs chemins de commande pouvaient créer une ambiguïté entre frontend, orchestrateur et exécuteur. L'architecture a été renforcée afin que le backend construise le manifeste autoritaire et que les agents ne fournissent pas de commandes shell libres.

\subsection{Compatibilité des runtimes Angular}

Lors des premiers essais Angular, le profil détecté Node.js 22.23.1 et npm 10.9.8 était rejeté par un catalogue ne reconnaissant qu'un profil plus ancien pour le premier stage. La résolution a consisté à ajouter un catalogue versionné contenant les paires validées tout en conservant l'ancienne version pour la reproductibilité.

\subsection{Persistance et migrations de base}

L'évolution rapide du backend a produit plusieurs migrations SQLite et, dans un cas, plusieurs têtes Alembic. Le diagnostic a nécessité l'identification de la base réellement utilisée et la réconciliation des migrations avant le redémarrage du workflow.

\subsection{Intégration du frontend}

L'intégration d'un nouveau cockpit avec des fonctionnalités existantes, notamment l'assistant et les points de décision, a demandé une fusion contrôlée. Les vérifications ont distingué les échecs préexistants des régressions introduites et ont utilisé une base temporaire pour les essais runtime.

\begin{table}[H]
    \centering
    \small
    \arrayrulecolor{cgigray}
    \rowcolors{2}{cgilight!55}{white}
    \begin{tabularx}{\textwidth}{|>{\bfseries\RaggedRight\arraybackslash}p{3.8cm}|>{\RaggedRight\arraybackslash}p{4.6cm}|>{\RaggedRight\arraybackslash}X|}
        \hline
        \rowcolor{cgilight!95}
        \textbf{Difficulté} & \textbf{Cause principale} & \textbf{Correction} \\
        \hline
        Patch non applicable & CRLF, contexte ou fichier obsolète. & Normalisation, checksum et validation préalable. \\
        \hline
        Commande ambiguë & Plusieurs composants pouvaient influencer l'exécution. & Manifeste construit par le backend. \\
        \hline
        Profil Angular bloqué & Catalogue incompatible avec le runtime validé. & Catalogue versionné et paire Node/npm explicite. \\
        \hline
        État de base incohérent & Migrations multiples ou base non autoritaire. & Réconciliation et identification de la base active. \\
        \hline
        Régression frontend & Fusion d'interfaces et fonctionnalités concurrentes. & Intégration par étapes et comparaison avec le baseline. \\
        \hline
    \end{tabularx}
    \caption{Synthèse des difficultés et corrections}
    \label{tab:difficultes-corrections-ch4}
    \rowcolors{2}{white}{white}
    \arrayrulecolor{black}
\end{table}

% ============================================================
\section{Fonctionnalités effectivement développées}
% ============================================================

\begin{table}[H]
    \centering
    \scriptsize
    \arrayrulecolor{cgigray}
    \rowcolors{2}{cgilight!55}{white}
    \begin{tabularx}{\textwidth}{|>{\bfseries\RaggedRight\arraybackslash}p{4.2cm}|>{\centering\arraybackslash}p{3.0cm}|>{\centering\arraybackslash}p{3.0cm}|>{\RaggedRight\arraybackslash}X|}
        \hline
        \rowcolor{cgilight!95}
        \textbf{Fonctionnalité} & \textbf{Java/Spring Boot} & \textbf{Angular} & \textbf{Observation} \\
        \hline
        Configuration et qualification & Implémentée & En intégration & Profils propres à l'écosystème. \\
        \hline
        Analyse et reviewer & Implémentés & Implémentés & Intégration Angular en cours. \\
        \hline
        Planification et reviewer & Implémentés & Implémentés & Trajectoires progressives. \\
        \hline
        Transformation isolée & Implémentée & Implémentée & Validation Angular globale en cours. \\
        \hline
        Build et tests & Implémentés & En validation & Maven contre npm/Angular CLI. \\
        \hline
        Réparation gouvernée & Implémentée & Planifiée & Modèle proposé/révisé/humain. \\
        \hline
        Cockpit et événements & Avancés & En intégration & Nouveau frontend Angular. \\
        \hline
        Assistant & Implémenté & À consolider & Lecture de l'état et des preuves. \\
        \hline
        Rapport final & Implémenté & En développement & Relié aux artefacts. \\
        \hline
        Authentification/RBAC & Non implémentés & Non implémentés & Évolution d'industrialisation. \\
        \hline
    \end{tabularx}
    \caption{État des fonctionnalités réalisées}
    \label{tab:fonctionnalites-realisees-ch4}
    \rowcolors{2}{white}{white}
    \arrayrulecolor{black}
\end{table}

\section*{Conclusion}
\addcontentsline{toc}{section}{Conclusion}

Ce chapitre a présenté l'implémentation de la Migration Factory à travers ses technologies, la structure des dépôts, les agents, l'orchestrateur, la persistance et le cockpit. La solution Java/Spring Boot dispose du parcours le plus complet et a permis de stabiliser les mécanismes de gouvernance, de réparation et de reporting. La solution Angular reprend les mêmes principes avec des profils, outils et agents adaptés, mais demeure en phase d'intégration et de validation globale.

L'implémentation montre que l'intelligence est volontairement séparée de l'autorité. Les agents analysent, planifient, transforment ou proposent ; le backend contrôle les transitions, les commandes, les chemins et l'application des correctifs. Cette séparation permet de profiter des capacités des LLM tout en conservant des preuves techniques et une supervision humaine.

Les prochaines étapes portent sur la consolidation du workflow Angular, l'enrichissement de l'évaluation quantitative, l'introduction d'une authentification et l'adaptation de la persistance à une exploitation multi-utilisateur.
